\chapter*{Abstract}\addcontentsline{toc}{chapter}{Abstract}
\label{chap:abstract}

High-dimensional proteomics can identify molecular differences between tumour and comparison tissue, but small cohorts, missing measurements, correlated proteins, incomplete laboratory metadata and information leakage create a substantial risk of irreproducible results. This thesis presents a from-scratch, data-only analysis of a supplied oral-cancer tissue proteomics matrix containing 84 specimens from 42 patients, each represented by a tumour and matched non-tumour specimen. The work integrates auditable data engineering, missingness-aware quality control, paired statistical inference, pathway analysis, patient-level heterogeneity assessment, leakage-controlled machine learning, multiverse sensitivity analysis and a locked prospective validation hand-off. Empirical analyses use only the supplied workbook. Research published by 31 December 2023 is used to justify methods but not to invent missing patient, laboratory or clinical information.

The workbook contains 8,071 supplied feature rows, 437,584 positive finite observations and 240,380 blank cells, corresponding to 35.456\% missingness. All 42 patient pairs are structurally complete, although only 1,350 feature rows are quantitatively complete across every pair. Quality control identified strong coupling between detection coverage and observed abundance, a broad tumour-associated multivariate structure, and seven patient pairs requiring sensitivity analysis rather than deletion. The primary no-imputation paired analysis identified 614 stable abundance candidates and 1,827 stable detection-pattern candidates under frozen effect, multiplicity and sensitivity rules. Reactome release 86 analysis found 271 pathways supported by both ranked enrichment and patient-level paired scoring. Positive tumour-associated organisation included RNA processing, translation and immune/antigen-handling sets; negative organisation included mitochondrial energy metabolism, extracellular-matrix degradation, complement and haemostasis. Consensus clustering did not support stable patient subtypes.

Nine predictive pipelines were evaluated using 25 repeated patient-grouped nested validations. The leading abundance elastic-net model achieved median ROC AUC 0.982, although a fold-local PCA baseline achieved 0.969 and the tuned elastic nets remained dense, preventing a compact-panel claim. A complete paired-label permutation pipeline produced an empirical p-value of 0.00498. Phase 6 stress testing evaluated 24 prespecified inferential scenarios and 42 exact leave-one-patient-out refits. Of the Phase 3 candidates, 583/614 abundance candidates and 1,804/1,827 detection candidates retained core support in at least 80\% of patient omissions; 596/783 Reactome pathways met joint direction and leading-edge robustness criteria. Frozen integration produced 165 high-priority internal candidates. Phase 7 converted these into a prospective, redundancy-aware hand-off of 12 abundance anchors and eight detection sentinels using branch-specific Pareto fronts and paired-change correlation components.

The principal conclusion is deliberately limited: the supplied paired-tissue matrix contains a strong and internally reproducible tumour-versus-matched-non-tumour signal across inferential, pathway and predictive analyses. It does not establish an externally validated biomarker, clinical diagnostic panel, causal mechanism, screening performance or clinical utility. The Phase 2 preprocessing benchmark remains unimplemented, and no independent cohort or orthogonal assay is available. These gaps are preserved as open dependencies rather than filled by unsupported assumptions.

\textbf{Keywords:} oral cancer; proteomics; paired design; missing data; empirical Bayes; pathway enrichment; elastic net; nested cross-validation; multiverse analysis; biomarker discovery; reproducibility.
